
\begin{result}{E.2.21}{The exact Artin image of the full plain integral CM diagram}{Proved}
Let \(I_{\mathrm{CM}}\) be the Artin image for the concrete
geometric integral diagram generated by all CM abelian varieties
over \(\Qb\), their actual isogeny lattices, specified tensor
and summand lattices, labelled algebraic correspondences and
all finite-model Artin algebra/module maps. The same assertion
holds with the labelled absolute-Hodge tensors of the abelian
version. Then
\[
 I_{\mathrm{CM}}
   =\Gamma_{F_{\mathrm{cm}}}\sqcup
        c\Gamma_{F_{\mathrm{cm}}},
 \qquad
 F_{\mathrm{cm}}=\prod_{K\ {\rm CM}}K^{\mathrm{ab}},
 \tag{E.2.21a}
\]
where the product denotes the compositum in \(\Qb\).
Its Artin kernel consists of the identity and the global
weight sign; each displayed fibre has exactly two elements.
In particular every \(\sigma\in\Sigma\) fixes the full
abelian extension of every imaginary-quadratic field,
including all its ray class fields.

The all-abelian criterion is consequently
\[
 \begin{split}
 \Sigma_F
  &=\{\sigma\in\Gamma_{F_{\mathrm{cm}}}:
       \rho_A(\sigma)(\Lambda_A)=\Lambda_{\sigma A}
       \text{ for every abelian variety }A/\Qb\},\\
 I_{\mathrm{ab}}&=\Sigma_F\sqcup c\Sigma_F .
 \end{split}
 \tag{E.2.21b}
\]
This identifies the CM restriction of every remaining
abelian Artin candidate without identifying a field
generated by CM division points with \(K^{\mathrm{ab}}\).
\end{result}
\begin{proof}
The CM tests used in E.2.5 and E.2.15 already belong
to the present CM/Artin diagram. They force the
restriction of an Artin candidate to
\(\Q^{\mathrm{cm}}\) to be the identity or actual
conjugation. Compose with conjugation and assume
\(\sigma\in\Gamma_{\Q^{\mathrm{cm}}}\).
The all-cyclic-quotient argument of E.2.9 also applies
inside this diagram, because a quotient of a CM
abelian variety is CM. Products give one common
scalar, and the Gaussian anchor makes its square
one. The two-type equations of E.2.20 are therefore
available for every Galois CM field \(F(i)\).

First, the imaginary-quadratic conclusion follows
directly from E.2.20. For a finite \(K\)-idele \(s\)
representing \(\sigma|_{K^{\mathrm{ab}}}\), the
identity CM type gives
\(\lambda s\in K^\times\). Since \(\lambda\) is now
one global sign, \(s\) is principal. Its finite
Artin image is trivial: the missing infinite
component belongs to \(\C^\times\), whose image
is trivial by connectedness. This proves fixation
of \(K^{\mathrm{ab}}\) through the actual reciprocity
map.

For the general assertion, let \(F/\Q\) be finite,
totally real and Galois, and put \(L=F(i)\).
For a finite \(L\)-idele \(s\) representing
\(\sigma|_{L^{\mathrm{ab}}}\), the quotient of the
two CM-type equations proves
\[
 s=b\,t,\qquad b\in L^\times,\quad
 t\in\mathbb A_{F,f}^{\times}.
 \tag{E.2.21c}
\]
This is the explicit Hilbert-90 calculation in
E.2.20. It cancels the common scalar, so this
part of the proof does not require its globality.
The principal factor has trivial finite Artin
image because \(L\) is CM.

Fix an arbitrary finite abelian extension \(M/L\).
Define the continuous homomorphism
\[
 v_M:\Gamma_F^{\mathrm{ab}}
   \xrightarrow{\ \operatorname{Ver}_{F\to L}\ }
   \Gamma_L^{\mathrm{ab}}
   \xrightarrow{\ \operatorname{res}_M\ }
   \Gal(M/L).
 \tag{E.2.21d}
\]
The first arrow is the actual transfer for the
index-two subgroup \(\Gamma_L\subset\Gamma_F\).
It is related to extension \(j_F^L\) of finite
ideles by
\[
 \operatorname{art}_L(j_F^L(u))
  =\operatorname{Ver}_{F\to L}
       \bigl(\operatorname{art}_F(u)\bigr).
 \tag{E.2.21e}
\]
The finite Galois transfer and Artin square are
\cite[Chapter V, Proposition 3.18 and adjoining
diagram, p.~162]{GenMilneCFT2020}.
For a finite target, choose an ideal representative
prime to its modulus to pass from that square to
finite ideles; weak approximation and the
principal-idele relation justify the replacement.
Taking inverse limits gives (E.2.21e).

For each real embedding \(\tau:F\hookrightarrow\R\),
let \(c_\tau\in\Gamma_F\) be its actual complex
conjugation. It has order two and lies outside
\(\Gamma_L\), since \(i\) is nonreal at every real
place. Using the two coset representatives
\(1,c_\tau\) gives the exact formula
\[
 \operatorname{Ver}_{F\to L}(c_\tau)=c_\tau^2=1.
 \tag{E.2.21f}
\]
No Galois assumption on \(M/F\), or inversion
assumption on \(\Gal(M/L)\), enters this calculation.

The finite quotient \(v_M\) defines an abelian
extension \(N/F\): its fixing subgroup in
\(\Gamma_F\) is the inverse image of
\(\ker(v_M)\). Equation (E.2.21f) kills the
complex conjugation at every real place.
Consequently every embedding of \(N\) is real,
so \(N\) is totally real.

Choose \(F'/\Q\) to be the normal closure of
\(N/\Q\). It is totally real and contains \(F\).
Moreover \(F'/F\) is abelian: since \(F/\Q\)
is Galois, each conjugate of \(N\) is an abelian
extension of \(F\), and their compositum remains
abelian over \(F\). Put \(L'=F'(i)\).
Then \(L'/\Q\) is Galois and CM, and \(L'/L\)
is abelian. Apply (E.2.21c) at this larger
field to obtain
\[
 s'=b't',\qquad
 \operatorname{art}_{L'}(s')
      =\sigma|_{L'^{\mathrm{ab}}},
 \quad b'\in L'^\times,\quad
 t'\in\mathbb A_{F',f}^{\times}.
\]

The norm/restriction Artin diagram gives
\[
 \begin{split}
 \sigma|_{L^{\mathrm{ab}}}
  &=\operatorname{art}_L
       \bigl(\operatorname{Nm}_{L'/L}(s')\bigr)\\
  &=\operatorname{art}_L
       \bigl(j_F^L(\operatorname{Nm}_{F'/F}(t'))\bigr).
 \end{split}
 \tag{E.2.21g}
\]
For the first equality, the actual field map is
\(L^{\mathrm{ab}}L'\subset L'^{\mathrm{ab}}\),
obtained by abelian base change.
For the second, the principal norm of \(b'\)
has trivial Artin image in the totally imaginary
field \(L\). The remaining norm identity follows
from
\[
 F'\otimes_F L=L',
\]
because \(L=F(i)\) and \(F'\) is totally real:
the determinant defining the norm commutes with
this extension of scalars.

Set
\[
 a=\operatorname{art}_F
       \bigl(\operatorname{Nm}_{F'/F}(t')\bigr).
\]
The same norm/restriction diagram makes \(a\)
trivial on \(N\), since \(N\subset F'\) and
\(\operatorname{art}_{F'}(t')\) fixes \(F'\).
Thus \(a\in\ker(v_M)\). Combining (E.2.21d),
(E.2.21e), and (E.2.21g) yields
\[
 \sigma|_M=v_M(a)=1.
\]
As \(M/L\) was arbitrary finite abelian, \(\sigma\)
fixes \(L^{\mathrm{ab}}\) pointwise.

Every finite CM field \(K\) is contained in
a Galois CM field \(L\) containing \(i\), obtained
by adjoining \(i\) to its normal closure.
For \(F=L^+\), one has \(L=F(i)\), with \(F/\Q\)
Galois and totally real. The preceding argument
applies. Since
\[
 K^{\mathrm{ab}}L\subset L^{\mathrm{ab}},
\]
\(\sigma\) fixes \(K^{\mathrm{ab}}\).
Varying \(K\) proves
\(\sigma\in\Gamma_{F_{\mathrm{cm}}}\).

Conversely E.2.6 constructs an actual comparison
for every element of \(\Gamma_{F_{\mathrm{cm}}}\):
the main CM theorem with idele \(1\), the
Tate-integral quasi-isogeny argument, exterior
cohomology, all labelled correspondences and
projector lattices, and the finite-model
Artin descent give every required map.
Complex conjugation supplies the other coset.
This proves (E.2.21a).
The proof of the Artin kernel in E.2.15 applies
verbatim to CM quotients and products, giving
the two global weight choices and exactly
two lifts in each fibre.

Restriction from the full geometric abelian
diagram gives the necessary field constraint
in (E.2.21b), and E.2.16 proves its exact
lattice necessity and sufficiency.
The known quartic construction already has
swapping elements in \(\Gamma_{F_{\mathrm{cm}}}\).
Thus the remaining all-abelian lattice condition,
the quartic place-pair, and extension to all
pure motives remain unresolved.
\end{proof}
